\section{问题四模型的建立与求解}

问题四用于检验模型对关键参数变化的敏感性。本文不采用单纯定性讨论，而是在每个情景下重新执行需求预测、站点选址和服务定价流程。

\subsection{模型的建立}

\subsubsection{扰动情景}

设基准情景为 $\Omega_0$。根据题目要求，设置如下情景：
\begin{align}
\Omega_1 &: \quad g=8\%,\\
\Omega_2 &: \quad p_i=5.5\%,\;q_i=9.5\%,\\
\Omega_3 &: \quad O_j'=1.2O_j,\\
\Omega_4 &: \quad B'=140.
\end{align}
四类情景分别对应老人增长加快、失能转移加重、运营成本上升和建设预算提高。

\subsubsection{鲁棒性指标}

设基准情景建设站点集合为 $J_0$，情景 $l$ 下建设站点集合为 $J_l$。站点稳定率为
\begin{equation}
\eta_l=\frac{|J_0\cap J_l|}{|J_0\cup J_l|}.
\end{equation}
设基准情景目标值为 $Z_0$，情景 $l$ 下最优目标值为 $Z_l^*$，基准方案应用于情景 $l$ 的目标值为 $Z_l(J_0)$，后悔值为
\begin{equation}
Reg_l=Z_l^*-Z_l(J_0).
\end{equation}
满意度变化率为
\begin{equation}
\Delta S_l=\frac{\bar{S}_l-\bar{S}_0}{\bar{S}_0}.
\end{equation}

\subsection{模型的求解}

对每个情景，本文依次执行以下步骤：重新计算第 5 年老人结构；重新计算理论需求和消费约束后需求；将更新需求输入问题二模型重新选址；固定新站点方案后重新求解问题三定价模型；最后比较站点数量、位置、服务定价、补贴总额、覆盖率、满意度和利润率。

从解释逻辑看，老人增长率和失能转移概率变化主要通过需求侧影响容量满足率；固定管理成本上升主要通过成本侧影响利润率和定价；预算增加直接扩大可行解集合，通常对站点规模和容量满足率影响最大。因此，问题四能够识别模型最敏感的参数来源。
